\documentclass[10pt,a4paper]{amsart}
\usepackage[margin=19mm]{geometry}
\usepackage{amsmath,amssymb,mathtools}
\usepackage[scr=rsfs,cal=euler]{mathalfa}
\usepackage{xfrac}
\usepackage[unicode,hidelinks,bookmarks=false]{hyperref}
\newcommand{\ie}{i.e.\ }
\newcommand{\cf}{cf.\ }
\newcommand{\resp}{respectively\ }
\newcommand{\Subsection}{\subsection}
\providecommand{\nobreakdash}{\nobreak\hbox{-}}
\setlength{\emergencystretch}{2em}
\multlinegap0pt
\usepackage{tikz-cd}
\theoremstyle{plain}
\newtheorem{lemma}{Lemma}
\newtheorem{theorem}[lemma]{Theorem}
\newtheorem{proposition}[lemma]{Proposition}
\newtheorem{corollary}[lemma]{Corollary}
\theoremstyle{definition}
\newtheorem{definition}[lemma]{Definition}
\theoremstyle{remark}
\newtheorem{remark}[lemma]{Remark}
\title[Automorphisms of Hilbert schemes]{Automorphisms of punctual Hilbert schemes and symmetric powers of varieties: the descent lemma for equivariant automorphisms}
\author{Ashima Bansal, Supravat Sarkar, Shivam Vats}
\date{Source: arXiv:2508.18059v2 (CC BY 4.0)}
\begin{document}
\maketitle
\noindent\emph{Source and licence.} Automorphisms of punctual Hilbert schemes and symmetric powers of varieties. Version arXiv:2508.18059v2 (CC BY 4.0), distributed under the Creative Commons Attribution 4.0 International licence, \url{http://creativecommons.org/licenses/by/4.0/}, as recorded in the arXiv licence metadata for this exact version and shown on the abstract page \url{https://arxiv.org/abs/2508.18059v2}. Full licence notice: licenses/arxiv-2508.18059-CC-BY-4.0.txt.\par\smallskip

\noindent\textit{Editorial scope.} Counted unit: the article's Lemma with source label \texttt{open} (source0.tex lines 772--774), the descent statement for an equivariant automorphism of the m-fold product of a product of curves over the m-fold product of the base curve: the automorphism either descends to an automorphism of each partial product, or the number of factors is two and the automorphism factors through a prescribed product; delivered together with the article's complete original proof of it (source0.tex lines 777--889, one \texttt{proof} environment of 113 lines, adjacent to the statement). No mathematics is written, repaired or re-derived: statement and proof are the article's own text line for line, in the article's own notation (the curves and their products, the equivariant automorphism, the partial products and their projections, the diagonal condition, the translation and bundle automorphisms), and no retained line is substituted, re-flowed or omitted. Required context retained uncounted, in source order: the statement of the article's lemma on the automorphism group of a product of curves (lines 560--563), and the statement of the article's corollary on automorphisms without unit components (lines 643--645), both of which the counted proof invokes. The proofs of those context results are not delivered. Nothing else of the article is retained: its main theorems and its remaining lemmas are omitted. No citation command occurs in any retained line, so the article's bibliography is not delivered and the wrapper carries no bibliography entry, although the article ships one. Every cross-reference inside the retained spans resolves inside the delivered document. Editorial formatting only: an AMS \texttt{amsart} preamble that restates the environments the retained lines use, namely the article's declarations of Lemma, Theorem, Proposition, Corollary, Definition and Remark sharing one counter as in the article, and that adds no mathematical macro, because every control sequence used by the retained lines is supplied by standard packages or by the wrapper's own mathematical runtime. The article's own source declares the class \texttt{amsart} as well; no class or style file is shipped with this package. Numbering differs from the article, because the article declares its Lemma environment as the master counter and resets it in each section, and because only a subset of environments is retained: the retained environments print with the article's shared counter in their source order, so the two retained context results print as Lemma 1 and Corollary 2 and the counted lemma as Lemma 3. No picture environment and no graphical inclusion occurs in any retained line, and the package delivers no image asset. One bundled source file, \texttt{sources/183830/source0.tex}, accompanies the delivered item as the exact source of every retained span. Every retained span is recorded in \texttt{delivery-evidence.json} with method \texttt{exact\_tex}.
\begin{lemma}\label{autoproduct}
Let $r, a_{1},\ldots, a_{r}$ be positive integers, and let $Y_{1},\ldots, Y_{r}$ be pairwise non-isomorphic indecomposable normal projective varieties. Suppose that either $H^{1}(Y_{i}, \mathcal{O}_{Y_{i}}) = 0$ for all $i$, or that each $Y_i$ is smooth with $K_{Y_i}$ ample. Let $Y\, =\, \prod _{i} Y_{i}^{a_{i}}$. Then, 
$$\textnormal{Aut}(Y)\, \cong \,\prod_{i}(\textnormal{Aut}(Y_{i})^{a_{i}}\,\rtimes\, S_{a_{i}}).$$
\end{lemma}

\begin{corollary}\label{nounit}
Let $E$ be an elliptic curve and suppose $\phi\,\in\, \emph{Aut}(E^m)$ descends to $S^m E$ and satisfies $\phi|_{\Delta_E}\,=\,$ id. Then $\phi$ is a permutation of the factors.
\end{corollary}

\begin{lemma}\label{open}
Let $C, F$ be smooth curves, $F$ (but not necessarily $C$) projective. Let $X\,=\, F\,\times\, C$, $X\,\xlongrightarrow{p}\,C$ the projection. Let $m\,\geq\, 2$, $\phi$ an $S_m$-equivariant automorphism of $X^{m}$ over $C^{m}$, such that $\phi{|_{\Delta_{X}}}\,=\, id$. For $1\,\leq\,i\,\leq\, m$, let $W_{i}\,=\, C^{i-1}\,\times\, X\,\times\, C^{m-i}$, $\pi=p_{(i-1)}\times id \times p_{(m-i)}: X\,\to \, W_{i}$. Then either $\phi$ descends to an automorphism of $W_{i}$ for all $i$, or $m\,=\,2$ and $\phi_{1}$ factors through $X\,\times\, X\,\xlongrightarrow{p\times id}\, C\,\times\,X$, $\phi_{2}$ factors through $X\,\times\, X\,\xlongrightarrow{id\times p}\, X\,\times\, C$.
\end{lemma}

\begin{proof}
We have the commutative diagram:

\begin{center}
\begin{tikzcd}
 (F\,\times\, C)^{m}  \arrow[rr, "\phi"] \arrow[swap]{dr}{p_{(m)}} &  _{}  & (F\,\times\,C)^{m}\arrow{dl}{p_{(m)}} \\[5pt]
 & C^{m}
\end{tikzcd}.
\end{center}
So, we get a morphism 
\begin{equation}\label{map to aut}
C^{m}\,\xlongrightarrow{T}\, \textrm{Aut}\,F^{m}.
\end{equation}

We will write elements of $X$ as $(z,c)$, with $z\,\in\, F$, $c\,\in\, C.$
We deal with the two cases separately.

\underline{Case\,1}: $F\,=\, \mathbb{P}^{1}$ or $g(F)\,\geq\, 2$.

By Lemma \ref{autoproduct}, $\textrm{Aut}(F^{m})\,=\, (\textrm{Aut}F)^{m} \rtimes\, S_{m}$. As, $S_{m}$ is discrete, there is $\sigma\,\in\, S_{m}$, $\tau_{i,\underline{c}}\,\in\, \textrm{Aut}(F)$ for $\underline{c}\,\in\, C^{m}$, $1\,\leq\,i\,\leq\,m$ such that
\begin{equation*}
T(\underline{c})(\underline{z})\,=\, (\tau_{i,\underline{c}}(z_{\sigma(i)}))_{i},    
\end{equation*} for all $\underline{z}\,\in\, F^m.$
So, 
\begin{equation}\label{horr p1}
\phi((z_{i},c_{i})_{i})\,=\, (\tau_{i,\underline{c}}(z_{\sigma(i)}), c_{i})_{i}. 
\end{equation}
If $\sigma\,=\, 1$, then $\phi$ descends to $W_{1}\,\xlongrightarrow{\psi_{1}}\,W_{1}$ given by 

\begin{equation*}
\psi_{1}((z_{1},c_{1}), c_{2},\ldots, c_{m}) \,=\, \left((\tau_{1,\underline{c}}(z_{1}), c_{1}),c_{2},\ldots,c_{m}\right)   
\end{equation*}
for $z_{1}\,\in\,F$, $\underline{c}\,\in\,C^{m}$. Similarly, $\phi$ descends to each $W_{i}$. 

So, assume $\sigma\,\neq\,1$. As $\phi$ is $S_m$-equivariant, for every $\pi\,\in\,S_{m}$, we have $\phi(\underline{x}_{\pi})\,=\, \phi(\underline{x})_{\pi}$. By \eqref{horr p1}, if $x_{i}\,=\, (z_{i},c_{i})\,\in\,F\,\times\, C,$
\begin{equation*}
\phi(\underline{x}_{\pi})\,=\, (\tau_{i, \underline{c}_{\pi}}(z_{\pi\,\sigma(i)}), c_{\pi(i)})_{i}     
\end{equation*}
and
\begin{equation*}
\phi(\underline{x})_{\pi}\,=\, (\tau_{\pi(i),\underline{c}}(z_{\sigma \pi{(i)}}), c_{\pi(i)})_i  
\end{equation*}
for all $i,\underline{c}, \underline{z}$. So, 
\begin{equation}\label{terr p1}
\tau_{i, \underline{c}_{\pi}}(z_{\pi\sigma{(i)}})\,=\, \tau_{\pi(i),\underline{c}}(z_{\sigma\pi(i)})  
\end{equation}
for all $\pi\,\in\, S_{m}$, $\underline{c}\,\in\, C^{m}, \underline{z}\,\in\,F^{m}$, $1\,\leq\, i\,\leq\, m$. If $m\,\geq\,3$, then since $\mathcal{Z}(S_{m})\,=\, 1$ and $\sigma\,\neq\, 1$, we can choose $\pi\,\in\, S_{m}$ such that $\pi\,\sigma\,\neq\, \sigma\,\pi$. Choose $i$ such that $\pi\,\sigma(i)\,\neq\,\sigma\,\pi(i)$. As $\eqref{terr p1}$ is true for all $\underline{z}$, we see that $\tau_{i,\underline{c}_{\pi}}$, $\tau_{\pi(i),\underline{c}}$ has to be constant function. But they are in $\textrm{Aut}\,F$, a contraction. 

So, $m\,=\,2$ and $\sigma$ is the transposition. Define


\begin{equation*}
\psi_{1}\,:\, C\,\times \,X\,\longrightarrow\, X    
\end{equation*}
\begin{equation*}
(c_{1},(z_{2}, c_{2}))\,\longmapsto\, (\tau_{1,\underline{c}}(z_{2}),c_{1})    
\end{equation*}
\begin{equation*}
\psi_{2}\,:\, X\,\times\,C\,\longrightarrow\, X    
\end{equation*}
\begin{equation*}
((z_{1},c_{1}), c_{2})\,\longmapsto\, (\tau_{2,\underline{c}}(z_{1}), c_{2}).    
\end{equation*}
Using $\eqref{horr p1}$, one sees that $\phi_{i}$ factors through $\psi_{i}$.


\underline{Case\,2}: $F\,=\, E$, an elliptic curve.

We have $\textrm{Aut}(E^{m})\,=\, E^{m}\,\rtimes\,\textrm{End}E^{m}$. As $\textrm{End}(E^{m})$ is discrete, there is $P\,\in\, \textrm{End}(E^{m})$ and  $\underline{a}_{\underline{c}}\,\in\, E^{m}$ for $\underline{c}\,\in\, C^{m}$ such that 
\begin{equation*}
T(\underline{c})(\underline{z})\,=\, P(\underline{z})\,+\, \underline{a}_{\underline{c}}.  
\end{equation*}
Let $P_{i}\,:\,E^{m}\,\longrightarrow\, E$ be the components of $P$, $a_{\underline{c},i}\,\in\, E$ the components of $\underline{a}_{\underline{c}},$ $1\,\leq\,i\,\leq\,m$. So, 
\begin{equation}\label{horr E}
\phi((z_{i},c_{i})_{i})\,=\, (P_{i}(\underline{z})+a_{\underline{c},i},c_{i})_{i}.    
\end{equation}
If $P\,=\,{id}$, then $\phi$ descends to $W_{1}\,\xlongrightarrow{\psi_{1}}\,W_{1}$ given by $\psi_{1}((z_{1},c_{1}), c_{2},\ldots,c_{m})\,=\, ((z_{1}+a_{\underline{c},1},c_{1}),c_{2},\ldots,c_{m})$, for $z_{1}\,\in\, E, \underline{c}\,\in\, C^{m}$. Similarly, $\phi$ descends to each $W_{i}$. 

So, assume $P\,\neq\,{id}$. As $\phi$ is $S_m$-equivariant, for every $\sigma\,\in\,S_{m}$, we have 
\begin{equation*}
\phi(\underline{x}_{\sigma})\,=\, \phi(\underline{x})_{\sigma} \quad \textrm{for all}\,\, \underline{x}\,\in\, X^{m}.
\end{equation*}
By \eqref{horr E}, if $x_{i}\,=\, (z_{i},c_{i})\,\in\, E\,\times\,C$,
\begin{equation*}
\phi(\underline{x}_{\sigma})\,=\, (P_{i}(\underline{z}_{\sigma})\,+\, a_{\underline{c}_{\sigma},i}, c_{\sigma(i)})_{i},    
\end{equation*}
and
\begin{equation*}
\phi(\underline{x})_{\sigma}\,=\, (P_{\sigma(i)}(\underline{z})\,+\, a_{\underline{c},\sigma(i)}, c_{\sigma(i)} )_{i}.   
\end{equation*}
 So, we have
\begin{equation*}
P_{i}(\underline{z}_{\sigma})\,+\, a_{\underline{c}_{\sigma},i}\,=\, P_{\sigma(i)}(\underline{z})\,+\, a_{\underline{c}, \sigma(i)} 
\end{equation*}
for all $\sigma\,\in\, S_{m}$, $\underline{c}\,\in\, C^{m},$ $ \underline{z}\,\in\, E^{m}$, $1\,\leq\,i\,\leq\, m$. This forces
$P_{i}(\underline{z}_{\sigma})\,=\, P_{\sigma(i)}(\underline{z})$. Hence, $P\,:\, E^{m}\,\longrightarrow\, E^{m}$ is $S_{m}$-equivariant, and $P|_{\Delta_{E}}\,=\,{id}$ as $\phi{|_{\Delta_{E}}}\,=\,{id}$. By Corollary \ref{nounit}, there is $\tau\,\in\, S_{m}$ such that $P(\underline{z})\,=\, \underline{z}_{\tau}$.

As $P$ is $S_{m}$--equivariant, we get $\tau\,\in\, \mathcal{Z}(S_{m})$. As $P\,\neq\, {id}$, we get $\mathcal{Z}(S_{m})\neq 1$. So $m\,=\,2$ and $\tau$ is the transposition. Define
\begin{equation*}
\psi_{1}\,:\, C\,\times\, X\,\longrightarrow\, X    
\end{equation*}
\begin{equation*}
(c_{1}, (z_{2}, c_{2}))\,\mapsto \, (z_{2}\,+\, a_{\underline{c},1},c_{1})    
\end{equation*}
and 
\begin{equation*}
\psi_{2}\,:\, X\,\times\, C\,\longrightarrow\, X    
\end{equation*}
\begin{equation*}
((z_{1}, c_{1}), c_{2})\,\mapsto\,(z_{1}\,+\, a_{\underline{c},2}, c_{2}).    
\end{equation*}
Using $\eqref{horr E},$ one sees that $\phi_{i}$ factors through $\psi_{i}$.
\end{proof}

\end{document}
